\documentclass[11pt,a4paper]{article}

%========================
% PACKAGES
%========================
\usepackage[utf8]{inputenc}
\usepackage[T1]{fontenc}
\usepackage[french]{babel}
\usepackage[a4paper,margin=2cm]{geometry}
\usepackage{amsmath,amssymb,amsthm}
\usepackage{lmodern}
\usepackage{fancyhdr}
\usepackage{lastpage}

%========================
% MISE EN PAGE
%========================
\pagestyle{fancy}
\fancyhf{}
\fancyhead[L]{Terminale Spé Maths}
\fancyhead[C]{Corrigé détaillé -- DM 0}
\fancyhead[R]{Page \thepage/\pageref{LastPage}}

%========================
\begin{document}

\section*{Exercice 1}

On considère $f(x)=x+\frac{1}{x}-2\ln x$ sur $]0,+\infty[$.

\subsection*{Partie A}

\textbf{1. Limite en $0^+$}

On sait que :
\[
\frac{1}{x}\to +\infty \quad \text{et} \quad \ln x \to -\infty
\]

Donc :
\[
-2\ln x \to +\infty
\]

Ainsi :
\[
f(x)\to +\infty
\]

\textbf{Conclusion :} $\lim_{x\to0^+} f(x)=+\infty$.

\bigskip

\textbf{2. Limite en $+\infty$}

On a :
\[
x\to+\infty,\quad \frac{1}{x}\to0,\quad \ln x\to+\infty
\]

Donc :
\[
f(x)=x-2\ln x + o(1)
\]

Or $x$ domine $\ln x$, donc :
\[
f(x)\to +\infty
\]

\bigskip

\textbf{3. Dérivée}

\[
f'(x)=1-\frac{1}{x^2}-\frac{2}{x}
\]

On met au même dénominateur :
\[
f'(x)=\frac{x^2 -1 -2x}{x^2} = \frac{(x-1)^2}{x^2}
\]

\bigskip

\textbf{4. Variations}

Comme $x^2>0$, le signe dépend de $(x-1)^2\ge0$.

Donc :
\[
f'(x)\ge0 \quad \text{et nul seulement en } x=1
\]

Donc $f$ est croissante sur $]0,+\infty[$.

Minimum en $x=1$ :
\[
f(1)=1+1-0=2
\]

\bigskip

\textbf{5. Inégalité}

Comme $f(x)\ge2$ :
\[
x+\frac{1}{x}-2\ln x \ge 2
\]

Donc :
\[
x+\frac{1}{x} \ge 2\ln x +2
\]

\bigskip

\subsection*{Partie B}

\textbf{6. Nombre de solutions}

$f$ est continue et strictement croissante, avec :
\[
f(1)=2<3,\quad f\to+\infty \text{ aux bornes}
\]

Donc 2 solutions.

\bigskip

\textbf{7. Position}

Une solution dans $]0,1[$ et une dans $]1,+\infty[$.

\bigskip

\textbf{8. Encadrement}

(on donne résultat numérique)

\[
x_1 \approx 0{,}28 \quad ; \quad x_2 \approx 3{,}57
\]

\bigskip

\section*{Exercice 2}

\subsection*{Partie A}

\textbf{1.} $g(0)=\frac{3}{2}$

\bigskip

\textbf{2.}

\[
g(x)=2-\frac{1}{x+2}
\]

\bigskip

\textbf{3.}

\[
\lim_{x\to+\infty} g(x)=2
\]

\bigskip

\textbf{4.}

\[
g'(x)=\frac{1}{(x+2)^2}>0
\]

Donc croissante.

\bigskip

\textbf{6. Point fixe}

\[
g(x)=x \Rightarrow x=\frac{2x+3}{x+2}
\]

\[
x(x+2)=2x+3
\]

\[
x^2=3 \Rightarrow x=\sqrt{3}
\]

\bigskip

\subsection*{Partie B}

\textbf{7.}

\[
u_1=\frac{7}{4},\quad u_2=\frac{13}{7}
\]

\bigskip

\textbf{8.} Par récurrence $u_n\ge1$.

\bigskip

\textbf{10.}

\[
u_{n+1}-u_n=\frac{(1-u_n)(u_n+1)}{u_n+2}
\]

Comme $u_n\ge1$, on a $u_{n+1}-u_n\le0$.

Donc suite décroissante.

\bigskip

\textbf{12.} Suite décroissante et minorée $\Rightarrow$ convergente.

\bigskip

\textbf{13. Limite}

\[
\ell=g(\ell) \Rightarrow \ell=\sqrt{3}
\]

\bigskip

\section*{Exercice 3}

\[
h(x)=x\sqrt{x^2+1}
\]

\textbf{1.} $h(0)=0$

\bigskip

\textbf{2.} Fonction impaire.

\bigskip

\textbf{3.}

\[
h'(x)=\frac{2x^2+1}{\sqrt{x^2+1}}
\]

\bigskip

\textbf{4. Tangente en 1}

\[
h(1)=\sqrt{2}, \quad h'(1)=\frac{3}{\sqrt{2}}
\]

\[
y=\frac{3}{\sqrt{2}}(x-1)+\sqrt{2}
\]

\bigskip

\textbf{5. Tangente parallèle à $y=3x-2$}

\[
h'(x)=3 \Rightarrow \frac{2x^2+1}{\sqrt{x^2+1}}=3
\]

Résolution $\Rightarrow x=\pm\sqrt{\frac{5+\sqrt{17}}{2}}$

\bigskip

\textbf{6. Dérivabilité en 0}

Taux de variation :
\[
\frac{h(x)-0}{x}=\sqrt{x^2+1}\to1
\]

Donc dérivable.

\bigskip

\textbf{7. Tangente en 0}

\[
y=x
\]

\bigskip

\section*{Exercice 4}

\textbf{1.}

\[
P(U_A)=\frac{1}{2},\quad P(U_B)=\frac{1}{2}
\]

\bigskip

\textbf{2.}

\[
P(R|A)=\frac{3}{5},\quad P(R|B)=\frac{1}{5}
\]

\bigskip

\textbf{3.}

\[
P(R)=\frac{1}{2}\times\frac{3}{5}+\frac{1}{2}\times\frac{1}{5}=\frac{2}{5}
\]

\bigskip

\textbf{4.}

\[
P(A|R)=\frac{\frac{3}{10}}{\frac{2}{5}}=\frac{3}{4}
\]

\bigskip

\textbf{5. Loi binomiale}

\[
X \sim \mathcal{B}(3,\frac{2}{5})
\]

\[
P(X=2)=\binom{3}{2}\left(\frac{2}{5}\right)^2\left(\frac{3}{5}\right)=\frac{36}{125}
\]

\[
P(X\ge1)=1-\left(\frac{3}{5}\right)^3=\frac{98}{125}
\]

\bigskip

\section*{Bonus}

\[
x+\frac{1}{x}=2+2\ln x
\]

C’est exactement l’égalité atteinte au minimum de $f$, donc solution unique :

\[
x=1
\]

\end{document}